\section{Офлайн-препроцессинг и сжатие графа}

Для обеспечения микросекундной задержки при маршрутизации вычислительное ядро полностью освобождено от задач парсинга геометрии. Все ресурсоемкие операции вынесены в специализированный офлайн-контур подготовки графа.

Инверсия топологии и реберный граф.
Классический узловой граф (где перекрестки — узлы, дороги — ребра) непригоден для мезоскопической модели, так как при такой архитектуре стоимость маневра (например, запрет поворота) зависит от истории предыдущего шага, что требует внедрения условных переходов в горячем цикле алгоритма. 

В разработанном контуре применена инверсия — направленный реберный граф:
\begin{itemize}
    \item Физические участки дорог (сегменты) становятся узлами нового вычислительного графа.
    \item Маневры на перекрестках (переходы между сегментами) становятся ребрами.
\end{itemize}
Это позволяет жестко внедрить все геометрические штрафы (за торможение, прохождение дуги поворота и разгон) в статический вес связующего маневра. Для графа размерности Московской агломерации генерация связей на перекрестках создает плотную сеть из 1.5 млн возможных маневров.

Хранение в формате компактных разреженных строк и топологическая пересортировка.
Классическое объектное представление графа в куче вызывает массовые промахи кэша процессора из-за случайного доступа к памяти. Для преодоления этого барьера («стены памяти») вся топология сериализуется в формат компактного хранения разреженных строк (CSR). Граф описывается набором плоских массивов (структура массивов):
\begin{enumerate}
    \item массив смещений указателей строк — массив из N+1 элементов, хранящий индексы начала списка смежных маневров для каждого узла.
    \item массив индексов столбцов — массив из E элементов, хранящий идентификаторы целевых сегментов.
    \item массив статических весов — массив из E элементов, хранящий базовое время проезда свободного ребра.
\end{enumerate}

Схематичная структура хранения компактного представления графа в бинарном дампе памяти приведена на рис. \ref{fig:binary_dump}.

\begin{figure}[htbp]
    \centering
    \includesvg[width=0.95\linewidth,height=0.45\textheight,keepaspectratio,inkscapelatex=false]{assets/pdf/binary_dump.drawio.svg}
    \caption{Структура бинарного представления графа компактного формата в памяти}
    \label{fig:binary_dump}
\end{figure}

Для максимизации процента попаданий в кэш процессора первого и второго уровня на этапе сборки дампа применяется алгоритм топологической пересортировки графа. Рекурсивная бисекция графа переназначает идентификаторы узлов таким образом, чтобы топологически связные сегменты получали смежные адреса в памяти. При загрузке узла контроллером памяти (блоками по 64 байта) аппаратно выполняется предвыборка его фактических соседей.

Выравнивание структур данных и предрасчет функций задержки.
В то время как статика графа хранится в формате структуры массивов, динамическое состояние (трафик) меняется непрерывно. Если хранить параметры затора разрозненно, вычисление функции задержки потребует нескольких промахов кэша на один узел. Для решения этой проблемы применена архитектура массива структур (AoS).

Состояние каждого физического сегмента инкапсулируется в структуру состояния динамического узла, которая принудительно выровнена под размер кэш-линии первого уровня директивой выравнивания по границе 64 байт:
\begin{plaincode}{dynamic_node_state}{Структура динамического состояния узла с выравниванием по кэш-линии}
struct alignas(64) DynamicNodeState {
    uint32_t K_magic;                         // Целочисленный множитель (4 байта)
    std::atomic<uint16_t> mpr_penalty;        // Внешний штраф (2 байта)
    uint16_t _reserved;                       // Заполнитель (2 байта)
    std::atomic<uint16_t> volume_buckets[24]; // Корзинки прогнозов (48 байт)
    // Оставшиеся 8 байт (защита от ложного разделения данных)
};
\end{plaincode}

Важнейшим элементом оптимизации является предрассчитанный целочисленный множитель. Для исключения операций деления и возведения в произвольную степень в горячем цикле роутера, базовые константы функции задержки (базовая скорость, емкость сегмента) аналитически сжимаются в единое целое число на этапе сборки графа:
\begin{equation}
    K_{magic} = \left\lfloor \frac{T_{free} \cdot \left( \frac{V_{free}}{5} - 1 \right)}{C_{300}^2} \cdot 2^{20} \right\rfloor
\end{equation}

Умножение на $2^{20}$ позволяет в реальном времени использовать аппаратную арифметику с фиксированной точкой (битовые сдвиги вправо на 20 разрядов) вместо медленных вычислений с плавающей запятой, сокращая время расчета динамического штрафа до 3–4 тактов процессора.

Генерация координационных матриц опорных вершин.
В процессе предварительной обработки также выполняется выборка опорных ориентиров (маяков) для алгоритма динамического поиска по опорным вершинам. Для равномерного покрытия транспортной сети применяется стратегия максимального покрытия в сочетании с алгоритмом граничного минмакса. Офлайн-генератор проводит полные обходы графа алгоритмом Дейкстры, фиксируя 32 наиболее равноудаленных маяка, после чего вычисляет и упаковывает расстояния от каждого узла до всех 32 маяков в бинарный дамп графа. Это формирует базу для вычисления эвристической оценки за константное время на лету.

Конвейер построения вычислительного графа.
Сквозной процесс офлайн-подготовки картографических данных представляет собой строго направленный конвейер (построитель графа), реализующий последовательную цепочку шагов от парсинга сырого картографического файла до выгрузки оптимизированного бинарного дампа графа. 

Основные этапы работы конвейера подготовки графа:
\begin{enumerate}
    \item Импорт геоданных: модуль импорта считывает сжатый бинарный файл выгрузки геоданных, извлекая сырые вершины, ребра и отношения. Данные загружаются во временную реляционную базу данных с пространственным расширением.
    \item Пространственная фильтрация: выполняется удаление тупиковых ребер, пешеходных и специализированных путей, незамкнутых островных подграфов, нарушающих свойство сильной связности графа. Выполняется привязка координат точек пересечений.
    \item Реберная трансформация графа: реляционный граф конвертируется в направленный реберный граф. Дорожные сегменты преобразуются в узлы вычислительного графа, а физически возможные маневры на перекрестках (с учетом знаков приоритета и запретов поворотов) — в направленные ребра.
    \item Расчет весов и оптимизация кэша: производится расчет статического времени преодоления свободных сегментов T\_free. Узлы переупорядочиваются методом обратного алгоритма Катхилла–Макки (RCM) для обеспечения максимальной локальности в памяти.
    \item Выбор ориентиров: запускается алгоритм выборки опорных маяков методом максимального покрытия. Выполняется расчет матрицы расстояний от всех узлов до маяков.
    \item Формирование структуры графа и сериализация: топология графа пакуется в плоские бинарные массивы компактного формата, выравнивается по размеру кэш-линии (64 байта) и сериализуется в единый оптимизированный бинарный файл графа, готовый для загрузки в оперативную память вычислительного ядра через системный вызов отображения файлов в память за константное время.
\end{enumerate}

Схематичная структура хранения компактного представления графа в бинарном дампе памяти приведена выше на рис. \ref{fig:binary_dump}.
